\documentclass[11pt, letterpaper]{report}
\input{../../homework.tex}
\usepackage{titlepageBU}
\title{Discrete Mathematics}
\date{10/31/24}
\courseID{MA 293}
\professor{Dr. Borkovitz}
\courseSection{A1}
\begin{document}
\makeproblem
\begin{notation}
	The set of natural numbers $\mathbb{N}$ is the set of all integers greater than or equal to $0$. That is, $\mathbb{N}=\left\{ 0,1,2,\ldots \right\} $. The set of integers $\mathbb{Z}$ is the set of integers. That is, $\mathbb{Z}=\left\{ \ldots,-1,0,1,\ldots \right\} $. The set of positive integers, $\mathbb{Z}^+$, is the set of all integers strictly greater than $0$. That is, $\mathbb{Z}^+=\left\{ 1,2,\ldots \right\} $.
\end{notation}
\begin{theorem}\label{thm:1}
	Let $n\in\mathbb{Z}^+$ be a positive integer. Then
	\[
		2+4+6+\cdots+2n=n(n+1)
	.\]
\end{theorem}
\begin{proof}
	We will use proof by the method of induction. Define $\Phi (n)$ as the statement
	\[
		\sum_{i=1}^{n} 2i=n(n+1)
	.\]
	We wish to show $\Phi (n)$ is true for all positive integers $n\in\mathbb{Z}^+$. Notice that for $n=1$, we have
	\[
		\sum_{i=1}^{1} 2i=2=1\cdot 2=1(1+1)=n(n+1)
	.\]
	Therefore, the statement holds for the base case $\Phi (1)$.

	Now suppose the statement $\Phi (k)$ holds for some positive integer $k$. That is, suppose
	\[
		\sum_{i=1}^{k} 2i=k(k+1)
	.\]
	We wish to show $\Phi (k+1)$ is true. That is, we wish to show
	\[
		\sum_{i=1}^{k+1} 2i=(k+1)((k+1)+1)
	.\]
	Notice that
	\begin{align*}
		\sum_{i=1}^{k+1} 2i&=2 \sum_{i=1}^{k+1} i\\
				   &= 2\left( (k+1)+\sum_{i=1}^{k} i \right) \\
				   &=2(k+1)+2 \sum_{i=1}^{k} i\\
				   &=2(k+1)+\sum_{i=1}^{k} 2i
	.\end{align*}
	By our inductive hypothesis, we know
	\[
		\sum_{i=1}^{k} 2i=k(k+1)
	.\]
	Thus, we have
	\begin{align*}
		2(k+1)+\sum_{i=1}^{k} 2i&=2(k+1)+k(k+1)\\
					&=(2k+2)+\left( k^2+k \right) \\
					&=k^2+3k+2\\
					&=(k+1)(k+2)\\
					&=(k+1)((k+1)+1)
	.\end{align*}
	Therefore, $\Phi (k+1)$ is true. By the method of induction, $2+4+\cdots+2n=n(n+1)$ for all positive integers $n\in\mathbb{Z}^+$.
\end{proof}
\begin{theorem}\label{thm:2}
	A tile game is played by two players who take turns. The game starts with $n$ red and $n$ blue tiles on the table, for $n$ some positive integer. On a turn, a player can take as many tiles as they wish, provided they are all from the same color. The last player to take a tile wins. In this game, player two has a winning strategy.
\end{theorem}
\begin{proof}
	We will use proof by the method of strong mathematical induction. Define a board position as a pair $(a,b)\in\mathbb{Z}^+$. Define $\Phi (n)$ as the statement that player two has a winning strategy, given an initial position $(n,n)$. We wish to show $\Phi (n)$ holds for all positive integers $n\in\mathbb{Z}^+$.

	Let $n=1$. It follows that the initial position is $(1,1)$. Since each player can only take one color, and must take at least one tile, the first player has only one possible move. They must take a single tile from a single color. Without loss of generality, suppose they take from color 1. Thus, the board for player 2's turn is $(0,1)$. Player two can now take the only remaining tile, outputting the board position $(0,0)$. Since player two took the last remaining tile, they win the game. Therefore, the statement $\Phi (1)$ is true.

	Now suppose $\Phi (k)$ holds for all positive integers $i\leq k$. That is, suppose all board positions of the form $(i,i)$ are winning for player two, provided it is player 1's move on the position $(i,i)$. We wish to show the statement $\Phi (k+1)$ holds; that is, we wish to show if the starting position is $(k+1,k+1)$, then there exists a winning strategy for player two.

	Suppose without loss of generality that player one takes $\alpha $ tiles from color 1, for $0<\alpha\leq k+1 $ a positive integer. It follows that the board position for player 2's turn is $(k+1-\alpha ,k+1)$. Notice that since $\alpha \geq 1$, we know $k+1-\alpha \leq k+1-1=k$, and thus $k+1-\alpha \leq  k$. Since our inductive hypothesis was that $\Phi (i)$ holds for all $0<i\leq k$, and $0<k+1-\alpha \leq k$, we know $\Phi (k+1-\alpha )$ holds. That is, we know $(k+1-\alpha ,k+1-\alpha )$ is a winning position for player 2, provided it is player 1's move, due to our inductive hypothesis. Using this logic, it follows that if player 2 also takes $\alpha $ tiles, but this time from color 2, then the resulting position on player 1's turn is $(k+1-\alpha ,k+1-\alpha )$, which by our inductive hypothesis we know to be winning for player 2. Therefore, the statement $\Phi (k+1)$ is true. By the method of induction, player two has a winning strategy for all intitial positions $(n,n)$, where $n\in\mathbb{Z}^+$.
\end{proof}
\setcounter{theorem}{3}
\begin{theorem}\label{thm:4}
	If you have a pile of $n$ tiles, split it into two smaller piles of tiles, multiply those numbers, repeat this process until you reach piles of 1 tile, then add up all multiples you found, you will end up with the number
	\[
		\frac{n(n-1)}{2}
	\]
	for any positive integer $n$.
\end{theorem}

Before beginning the proof, I feel it's important to consider the logic that led to the solution. If we want to define a function $\phi  $ that outputs the final total given a starting number of tiles $n$, we must first show our function is well defined, meaning for any given $n$, there exists exactly one value of $\phi  (n)$. More intuitively, we must show splitting the pile in different ways does \emph{not} yield a different final total. The proof of this requires strong induction and is fairly mathematical. However, this lets us define a precise value. In the proof, we will show the base case $\phi  (1)=0$. If we notice that we can decompose a pile of $n$ tiles into two piles of $n-1$ tiles and a pile of $1$ tile, then we can consider the product of $n-1$ and $1$ \emph{plus the value of }$\phi  (n-1)+\phi  (1)$, since the possible decompositions of $n-1$ tiles do not change if there happen to be decompositions that happened previously. Since $\phi  (1)=0$ and $(n-1)\cdot 1=n-1$, we claim $\phi  (n)=\phi  (n-1)+(n-1)$. Notice this is just $T(n-1)$ for $T$ the triangle numbers, and thus we claim the closed form formula $\phi  (n)=\frac{n(n-1)}{2}$ holds for all positive integers $n$. This is the rationale behind what we decided to prove.
\begin{proof}
	We will use proof by the method of strong induction. Let $\Phi (n)$ be the statement that if we decompose some number $n$ of tiles into two different numbers of tiles, multiply those numbers, perform the same decomposition and multiplication on the remaining pies until we reach piles of 1, then add together all the products, then the resultant number, which we define as $\phi  (n)$, is the same every time. Moreover, $\phi  (n)=\frac{n(n-1)}{2}$. We will prove $\Phi (n)$ holds for all positive integers $n\in\mathbb{Z}^+$.

	First, we show the base case. Suppose $n=1$. Since 1 tile cannot be split into two piles of tiles, we have a sum of 0 terms. Thus, $\phi  (1)=0$. Additionally, notice that
	\[
		\frac{1(1-1)}{2}=0
	.\]
	Thus, $\Phi (1)$ holds.

	Now, suppose for some positive integer $k$ that for all $0<i\leq k$, $\Phi (i)$ holds. We wish to show $\Phi (k+1)$ holds; that is, we wish to show $\phi  (k+1)$ has precisely one possible value, namely
	\[
		\phi  (k+1)=\frac{(k+1)((k+1)-1)}{2}=\frac{k(k+1)}{2}
	.\]
	First, we show $\phi  (k+1)$ has precisely one value. Suppose there exist positive integers $a,b,c,d\in\mathbb{Z}^+$ such that $a+b=k+1$ and $c+d=k+1$. These integers do not necessarily need to be distinct. Notice that if we decompose the pile $k+1$ into $a$ and $b$ tiles, we have
	\[
		\phi  (k+1)=ab+\phi  (a)+\phi  (b)
	.\]
	By the same logic, if we decompose the pile into $c$ and $d$ tiles, we have
	\[
		\phi  (k+1)=cd+\phi  (c)+\phi  (d)
	.\]
	We must show these two decompositions are equivalent in the general case. Specifically, we must show that $a+b=c+d$ implies they are equivalent. Since $a,b,c,d> 0$, we know $\Phi (a),\Phi (b),\Phi (c),\Phi (d)$ hold, and thus
	\[
		\phi  (a)=\frac{a(a-1)}{2},\quad \phi  (b)=\frac{b(b-1)}{2},\quad \phi  (c)=\frac{c(c-1)}{2},\quad \phi  (d)=\frac{d(d-1)}{2}
	.\]
	It follows that
\begin{align*}
	&\qquad\;\;\; a+b=c+d\\
	&\implies (a+b)^2=(c+d)^2\\
	&\implies a^2+2ab+b^2=c^2+2cd+d^2\\
	&\implies \frac{a^2+2ab+b^2}{2}=\frac{c^2+2cd+d^2}{2}\\
	&\implies ab+\frac{a^2+b^2}{2}=cd+\frac{c^2+d^2}{2}\\
	&\implies ab+\frac{a^2-b^2}{2}-\frac{a+b}{2}=cd+\frac{c^2+d^2}{2}-\frac{c+d}{2}\\
	&\implies ab+\frac{a^2-a+b^2-b}{2}=cd+\frac{c^2-c+d^2-d}{2}\\
	&\implies ab+\frac{a^2-a}{2}\cdot \frac{b^2-b}{2}=cd+\frac{c^2-c}{2}\cdot \frac{d^2-d}{2}\\
	&\implies ab+\frac{a(a-1)}{2}=cd+\frac{c(c-1)}{2}\\
	&\implies ab+\phi  (a)+\phi  (b)=cd+\phi  (c)+\phi  (d)
.\end{align*}
Therefore, there is precisely one value of $\Phi (k+1)$. It remains to be shown that $\phi  (k+1)=\frac{k(k+1)}{2}$.

Notice that $k+1=k+1$, and thus by previous logic, $\phi (k+1)=(k\cdot 1)+\phi  (k)+\phi  (1)$. It follows that
\begin{align*}
	\phi  (k+1)&=\phi  (k)+k\\
		  &=\frac{k(k-1)}{2}+k\\
		  &=\frac{k^2-k}{2}+\frac{2k}{2}\\
		  &=\frac{k^2+k}{2}\\
		  &=\frac{k(k+1)}{2}
.\end{align*}
Therefore, $\phi  (k+1)=\frac{k(k+1)}{2}$, and $\Phi (k+1) $ holds. By the method of strong induction, we have shown $\Phi (n)$ holds for all positive integer $n$. That is, we have shown
\[
	\phi  (n)=\frac{n(n-1)}{2}
.\]
\end{proof}
\section*{Acknowledgements}
I worked alone.

Academic Integrity -- Please type your name to acknowledge that if
you collaborated on this problem, then it was only with classmates
from either section of CAS MA 293 (except to have a partner for the
game), that you only used allowable resources (notes from class, no
Internet searches, AI, etc.), and that you wrote this paper by
yourself: \textbf{Grant} \textbf{Talbert}
\end{document}
